% GENERATED FILE. Edit canonical lessons, then run npm run generate:books.
% canonical-source-hash: fnv1a64:209dd3f8ccee764f
% canonical-lessons: ML-C333-kasuvanti, ML-C333-karimpu, ML-C333-arival, ML-C333-pottu, ML-C333-helmerru

\chapter{Cashew nut, Sugarcane, Sickle, Bindi, Helmet}
\label{ch:ml-cashew-nut-sugarcane-sickle-bindi-helmet}
\hlchaptermodality{333}

\begin{chapteropening}
\textbf{By the end of this chapter:} \emph{I can say a cashew nut, sugarcane, a sickle, a bindi and a helmet.}

Complete the last lesson of chapter 333: I can say a cashew nut, sugarcane, a sickle, a bindi and a helmet.
\end{chapteropening}

\section[\texorpdfstring{\ml{കശുവണ്ടി} (kaśuvaṇṭi)}{കശുവണ്ടി (kaśuvaṇṭi)}]{\ml{കശുവണ്ടി} (kaśuvaṇṭi) — a cashew nut}

\label{lesson:ML-C333-kasuvanti}



\begin{quote}
Before the new one: say the Malayalam for ash, then the Malayalam for a sapling.
\end{quote}

\subsection*{You'll want to know: \ml{കശുവണ്ടി}}
\textbf{\ml{കശുവണ്ടി}} — \emph{kaśuvaṇṭi} — \textquotedblleft{}a cashew nut\textquotedblright{}.

\begin{grammarlens}[title={What you've built}]
One more word for this chapter.
\end{grammarlens}

\subsection*{Your turn}
\begin{itemize}
\raggedright
  \item \emph{Say it:} \emph{kaśuvaṇṭi}
  \item \emph{Say it:} \emph{kaśuvaṇṭi}, once more
  \item \emph{Say it:} say \emph{kaśuvaṇṭi}
  \item \emph{Recall:} say the Malayalam for ash, then the Malayalam for a sapling, then say \emph{kaśuvaṇṭi} again
\end{itemize}

\begin{tcolorbox}[breakable,colback=teal!4,colframe=teal!35!black,title={Before you move on}]
What does \ml{കശുവണ്ടി} mean? (\textquotedblleft{}A cashew nut\textquotedblright{}.) Say it once more.
\end{tcolorbox}
\section[\texorpdfstring{\ml{കരിമ്പ്} (karimpŭ)}{കരിമ്പ് (karimpŭ)}]{\ml{കരിമ്പ്} (karimpŭ) — sugarcane}

\label{lesson:ML-C333-karimpu}



\begin{quote}
Before the new one: say the Malayalam for a sapling, then the Malayalam for a cashew nut.
\end{quote}

\subsection*{You'll want to know: \ml{കരിമ്പ്}}
\textbf{\ml{കരിമ്പ്}} — \emph{karimpŭ} — \textquotedblleft{}sugarcane\textquotedblright{}.

\begin{grammarlens}[title={What you've built}]
One more word for this chapter.
\end{grammarlens}

\subsection*{Your turn}
\begin{itemize}
\raggedright
  \item \emph{Say it:} \emph{karimpŭ}
  \item \emph{Say it:} \emph{karimpŭ}, once more
  \item \emph{Say it:} say \emph{karimpŭ}
  \item \emph{Recall:} say the Malayalam for a sapling, then the Malayalam for a cashew nut, then say \emph{karimpŭ} again
\end{itemize}

\begin{tcolorbox}[breakable,colback=teal!4,colframe=teal!35!black,title={Before you move on}]
What does \ml{കരിമ്പ്} mean? (\textquotedblleft{}Sugarcane\textquotedblright{}.) Say it once more.
\end{tcolorbox}
\section[\texorpdfstring{\ml{അരിവാൾ} (arivāḷ)}{അരിവാൾ (arivāḷ)}]{\ml{അരിവാൾ} (arivāḷ) — a sickle}

\label{lesson:ML-C333-arival}



\begin{quote}
Before the new one: say the Malayalam for a cashew nut, then the Malayalam for sugarcane.
\end{quote}

\subsection*{You'll want to know: \ml{അരിവാൾ}}
\textbf{\ml{അരിവാൾ}} — \emph{arivāḷ} — \textquotedblleft{}a sickle\textquotedblright{}.

\begin{grammarlens}[title={What you've built}]
One more word for this chapter.
\end{grammarlens}

\subsection*{Your turn}
\begin{itemize}
\raggedright
  \item \emph{Say it:} \emph{arivāḷ}
  \item \emph{Say it:} \emph{arivāḷ}, once more
  \item \emph{Say it:} say \emph{arivāḷ}
  \item \emph{Recall:} say the Malayalam for a cashew nut, then the Malayalam for sugarcane, then say \emph{arivāḷ} again
\end{itemize}

\begin{tcolorbox}[breakable,colback=teal!4,colframe=teal!35!black,title={Before you move on}]
What does \ml{അരിവാൾ} mean? (\textquotedblleft{}A sickle\textquotedblright{}.) Say it once more.
\end{tcolorbox}
\section[\texorpdfstring{\ml{പൊട്ട്} (poṭṭŭ)}{പൊട്ട് (poṭṭŭ)}]{\ml{പൊട്ട്} (poṭṭŭ) — a bindi, a forehead dot}

\label{lesson:ML-C333-pottu}



\begin{quote}
Before the new one: say the Malayalam for sugarcane, then the Malayalam for a sickle.
\end{quote}

\subsection*{You'll want to know: \ml{പൊട്ട്}}
\textbf{\ml{പൊട്ട്}} — \emph{poṭṭŭ} — \textquotedblleft{}a bindi, a forehead dot\textquotedblright{}.

\begin{grammarlens}[title={What you've built}]
One more word for this chapter.
\end{grammarlens}

\subsection*{Your turn}
\begin{itemize}
\raggedright
  \item \emph{Say it:} \emph{poṭṭŭ}
  \item \emph{Say it:} \emph{poṭṭŭ}, once more
  \item \emph{Say it:} say \emph{poṭṭŭ}
  \item \emph{Recall:} say the Malayalam for sugarcane, then the Malayalam for a sickle, then say \emph{poṭṭŭ} again
\end{itemize}

\begin{tcolorbox}[breakable,colback=teal!4,colframe=teal!35!black,title={Before you move on}]
What does \ml{പൊട്ട്} mean? (\textquotedblleft{}A bindi, a forehead dot\textquotedblright{}.) Say it once more.
\end{tcolorbox}
\section[\texorpdfstring{\ml{ഹെൽമെറ്റ്} (helmeṟṟŭ)}{ഹെൽമെറ്റ് (helmeṟṟŭ)}]{\ml{ഹെൽമെറ്റ്} (helmeṟṟŭ) — a helmet}

\label{lesson:ML-C333-helmerru}



\begin{quote}
Before the new one: say the Malayalam for a sickle, then the Malayalam for a bindi.
\end{quote}

\subsection*{You'll want to know: \ml{ഹെൽമെറ്റ്}}
\textbf{\ml{ഹെൽമെറ്റ്}} — \emph{helmeṟṟŭ} — \textquotedblleft{}a helmet\textquotedblright{}.

\begin{grammarlens}[title={What you've built}]
One more word for this chapter.
\end{grammarlens}

\subsection*{Your turn}
\begin{itemize}
\raggedright
  \item \emph{Say it:} \emph{helmeṟṟŭ}
  \item \emph{Say it:} \emph{helmeṟṟŭ}, once more
  \item \emph{Say it:} say \emph{helmeṟṟŭ}
  \item \emph{Recall:} say the Malayalam for a sickle, then the Malayalam for a bindi, then say \emph{helmeṟṟŭ} again
\end{itemize}

\begin{tcolorbox}[breakable,colback=teal!4,colframe=teal!35!black,title={Before you move on}]
What does \ml{ഹെൽമെറ്റ്} mean? (\textquotedblleft{}A helmet\textquotedblright{}.) Say it once more.
\end{tcolorbox}
